% ---------------------------------------------------------------
% Appendix C body text: "Additional experiments"
%
% One block per subsection. Each block is the prose only - paste it
% under the \subsection heading that already exists in the draft; the
% headings are repeated here as comments so the blocks are easy to place.
%
% Figure labels assumed (from paper_final_artifacts/figures.tex):
%   \label{fig:activation-geometry}   Figure 1
%   \label{fig:methods}               Figure 2
%   \label{fig:pooled-training}       Figure 4
%   \label{fig:heldout-targets}       Figure 5
%   \label{fig:supp-framing}          Figure A1
%   \label{fig:supp-depth}            Figure A2
%   \label{fig:supp-features}         Figure A3
%   \label{fig:supp-label-types}      Figure A4
%   \label{fig:supp-geometry}         Figure A5
% ---------------------------------------------------------------


% ===============================================================
% C.1  Unsupervised characterization of True/False prompt framing
%      versus original prompts
% ===============================================================

Figure~\ref{fig:supp-framing} shows that the prompt template, rather than the
binary label, accounts for the largest share of variation once both prompt
formats are placed on a common map. The PCA behind this figure is fit on the
original and True/False activations pooled, so that the framed prompts help
define the basis they are drawn in; projecting them into a basis fit on the
original prompts alone would make their distance from it partly an artifact of
the fit. The two framings occupy largely disjoint regions in both models and
separate along the leading component (24.7\% of the variance for Gemma, 18.3\%
for Qwen), whereas the two label classes overlap throughout. Measured on the
first four components, the two framings' mean activations are separated by
1.1--1.2 times the root-mean-square spread of the full cloud, against 0.4--0.5
for the two label classes. Panel (a) shows that the dataset structure of
Figure~\ref{fig:activation-geometry} is reproduced inside each framing region:
most datasets appear as a pair of clusters, one per template. Social IQa is the
exception, as it has no original-prompt variant, so its two framings are
identical activations and it is omitted from panel (b).

This geometry constrains how framing can be used. Probes are fit and evaluated
within a single framing throughout the paper, since a probe fit on one template
would otherwise be evaluated in a region of activation space that its training
data does not cover. The displacement does not remove label information:
converting heterogeneous prompts into one declarative template improves
transfer to held-out targets by up to 0.22 AUC
(Figure~\ref{fig:heldout-targets}) and raises out-of-domain performance at
every depth past 53\% (Figure~\ref{fig:supp-depth}). Prompt framing is
therefore better treated as a change of coordinates that probes must be refit
for than as a perturbation that a fitted probe can be expected to tolerate.


% ===============================================================
% C.2  Analyzing probe performance across different layers
% ===============================================================

Figure~\ref{fig:supp-depth} shows that probe performance is stable from 68\%
depth onward, and that the depth at which transfer becomes possible depends on
the model. Within-target performance is nearly flat: the median best-method AUC
across datasets varies by at most 0.02 between 53\% and 100\% depth for every
model and framing, so the choice of layer is close to immaterial when the
target dataset is included in training. Under transfer the earliest depth
behaves differently in the two models. Qwen at 53\% depth trails its own best
later layer by 0.16 AUC when the target's domain is absent and by 0.20 AUC
under one-source training, whereas Gemma trails by at most 0.08 and is
marginally best at 53\% under domain-absent training on original prompts. From
68\% onward every curve is flat to within 0.03 between consecutive depths, and
True/False framing stays above original framing in both models at every depth
in both transfer panels.

The disagreement at 53\% depth is the reason we replicate each layer-dependent
claim across both models and both framings rather than reporting a single
layer. A study reading Qwen at an early layer would conclude that out-of-domain
transfer is weak, and a study reading Gemma at the same layer would conclude
that it is already near its ceiling; neither conclusion is a property of linear
probing. Because the curves converge above 68\%, the main figures either fix
88\% depth or aggregate over the four depths, and in that band the choice of
layer contributes less to probe performance than the training regime does
(Figure~\ref{fig:methods}).


% ===============================================================
% C.3  Comparing raw activations to alternative activation-based
%      features
% ===============================================================
% NOTE: the draft already carries the SAE motivation sentence
% ("Sparse autoencoders (SAEs) have achieved widespread importance ...
% we were consequently interested in how these features performed in
% probing."). The text below is written to follow it.

We compare four alternatives to the raw last-token activation, each fit on
training folds only: PCA retaining 95\% of the training variance, the per-layer
SAE and transcoder dictionary activations, the three concatenated (PCA +
dictionary), and a 100-bin histogram of the values in each activation vector.
Dictionary features exist for Gemma only, since pretrained per-layer
dictionaries are available for that model and not for Qwen (Appendix A.5).
Figure~\ref{fig:supp-features} shows that none of these representations
improves on the raw activation. Across 68--100\% depth, PCA at 95\% variance is
effectively lossless, with a median change of 0.000 AUC under one-source
training and $-0.004$ within target; SAE, transcoder, and PCA + dictionary
features lose 0.01--0.04 AUC at the median. Sparse dictionary features are
therefore close enough to the raw activation to be usable for probing, but
nothing in the learned basis recovers separability that the raw vector lacks.

Panel (a) uses the value histogram as a control on what neuron identity is
worth: it preserves the marginal distribution of activation values and discards
which coordinate produced each value. The histogram loses 0.004 AUC on Geometry
of Truth and 0.006 on Harmful/Harmless, so those datasets are separable from
the activation marginal alone, and the loss grows to 0.13 on VDS, 0.16 on
WJB-Adversarial, and 0.31 on Anatomy; the median loss over all datasets is
about 0.10. The datasets that depend most on coordinate-specific structure are
the same ones that require dataset-specific random slopes in
Figure~\ref{fig:pooled-training}, which is consistent with their separating
directions lying outside the subspace the remaining datasets share.


% ===============================================================
% C.4  Assessing probe transfer across related concepts
% ===============================================================

Figure~\ref{fig:supp-label-types} shows that class-mean directions transfer
between datasets whose labels assert different things. We partition the
datasets by what $y=1$ means: True/False (the displayed statement is true),
Yes/No (the correct answer to the displayed question is yes), and Safe/Unsafe
(the displayed request is safe). Each cell is the median diff-in-means AUC over
the ordered source--target dataset pairs of two label types at 68\% and 88\%
depth on original prompts, with the ten pooled datasets as sources and all
twelve as targets. Every cross-type cell is above chance, at 0.61--0.74 for
Gemma and 0.77--0.86 for Qwen, so the three label conventions are sign-aligned:
a direction fit to separate true statements from false ones scores safe
requests above unsafe ones with no relabeling.
Within-type transfer is the largest value in its row in four of the six rows;
in the two True/False rows it ties with transfer to Yes/No targets (0.800
against 0.802 for Gemma, 0.857 against 0.857 for Qwen), which is expected given
that the correct answer to a Yes/No question is itself a truth claim. The Gemma
matrix is asymmetric: directions fit on the safety datasets reach True/False
and Yes/No targets at 0.66 and 0.61, while directions fit on those two types
reach the safety datasets at 0.73 and 0.74.

This partition is not a relabeling of the domain structure used elsewhere in
the paper. The medical domain splits across it, with Anatomy in the True/False
group and the two PubMedQA sets forming the Yes/No group, so the domain-absent
regime and these label types withhold different training data. Read together
with the higher within-type values, the cross-type results indicate that truth
and safety are read out along related but not interchangeable directions, and
that the relationship is stronger in Qwen than in Gemma. The Safe/Unsafe cells
should be read alongside the fact that the safety datasets are also the most
separable within target (Figure~\ref{fig:pooled-training}), which raises both
the Safe/Unsafe diagonal and the columns leading into it.


% ===============================================================
% C.5  Predicting transfer from label-free activation geometry
% ===============================================================

Figure~\ref{fig:supp-geometry} shows that activation geometry predicts transfer
between datasets in the same domain but not between domains. For each of the 16
model $\times$ framing $\times$ depth settings we rank the 110 ordered
source--target pairs (ten sources, twelve targets) by three similarity measures
and correlate each ranking with the AUC of an $L_2$ probe trained on the source
and evaluated on the target. The probe-direction cosine, the regularized
Mahalanobis cosine of Appendix B.2, requires the target's labels and ranks
transfer almost perfectly, at 0.87 for same-domain pairs and 0.96 for
different-domain pairs. The two label-free measures separate sharply. The
overlap of the datasets' 20 leading principal components ranks transfer at 0.59
within a domain and 0.05 across domains, and the whitened distance between the
datasets' mean activations ranks it at 0.12 within a domain and 0.00 across
domains. Second-moment structure therefore carries considerably more
information about transfer than the distance between dataset means does, and
neither measure is informative once the two datasets belong to different
domains.

Panel (b) shows the mechanism behind the within-domain result. For pairs in the
same domain, the absolute cosine between the two datasets' probe directions
falls from 0.56 to 0.28 as the distance between their activation centroids
grows, while for pairs in different domains it stays at 0.23--0.29 and does not
depend on that distance. Proximity in activation space tracks agreement between
separating directions only after the two datasets already share a domain. This
bounds one construction used elsewhere in the paper: when a held-out target has
no other dataset in its domain, the random-slopes prediction weights the
available dataset adjustments by exactly this whitened centroid distance
(Appendix B.3), which is the weakest of the three measures evaluated here, and
the VDS rows of Figure~\ref{fig:heldout-targets} should be read with that
limitation in mind. Selecting a transferable source dataset without target
labels remains unresolved, and the covariance subspace measure is the more
promising of the two label-free candidates.
